\chapter{Kết luận}
\label{chap:conclusion}

Chương này tổng kết các kết quả nghiên cứu của luận văn và đánh giá ý nghĩa của phương pháp PhaBERT được đề xuất trong Chương~\ref{chap:method} và kiểm chứng tại Chương~\ref{chap:thuc_nghiem_danh_gia}. Nội dung được tổ chức theo ba phần: tóm tắt kết quả và nhận định chính, tổng hợp các đóng góp của nghiên cứu, và thảo luận về hạn chế cùng hướng phát triển trong tương lai.

\section{Kết luận}
\label{sec:conclusion}

Luận văn này nghiên cứu bài toán phân loại hệ gen thực khuẩn thể từ dữ liệu hệ gen môi trường, một thách thức quan trọng trong sinh tin học hiện đại. Các phương pháp hiện có gặp khó khăn khi xử lý contig ngắn do phụ thuộc vào công cụ dự đoán gen như Prodigal~\cite{hyatt2010prodigal}, vốn không hoạt động hiệu quả trên các đoạn DNA ngắn thiếu khung đọc mở hoàn chỉnh~\cite{trimble2012short}. Để giải quyết vấn đề này, luận văn đề xuất PhaBERT, một kiến trúc học sâu lai kết hợp mô hình nền tảng hệ gen DNABERT-2~\cite{zhou2023dnabert} với mô-đun MKC gồm nhánh CNN đa cửa sổ và gộp trung bình có mặt nạ toàn cục. Phương pháp hoạt động trực tiếp trên chuỗi DNA thô thông qua mã hóa BPE, không phụ thuộc vào công cụ dự đoán gen, và được huấn luyện theo chiến lược hai giai đoạn với tốc độ học phân biệt để tối ưu hóa đồng thời mạng nền và các thành phần chuyên biệt cho tác vụ.

Phương pháp được đánh giá toàn diện trên bốn nhóm độ dài contig (100--400 bp, 400--800 bp, 800--1.200 bp, 1.200--1.800 bp) thông qua kiểm định chéo phân tầng 5 phần trên tập dữ liệu gồm 2.241 hệ gen thực khuẩn thể hoàn chỉnh. Như trình bày trong Bảng~\ref{tab:main_results}, PhaBERT đạt độ chính xác cao nhất trên tất cả bốn nhóm contig với các giá trị lần lượt là 82,01\%, 87,33\%, 89,90\% và 91,38\%. So với các phương pháp tiên tiến, PhaBERT cải thiện 0,36--3,09\% so với PhaTYP~\cite{shang2023phatyp}, 2,29--3,94\% so với DeePhage~\cite{wu2021deephage}, và 1,33--6,65\% so với ProkBERT~\cite{ligeti2024prokbert}.

Kết quả nghiên cứu chứng minh tính hiệu quả của việc kết hợp mô hình nền tảng hệ gen tiền huấn luyện với kiến trúc chuyên biệt MKC cho bài toán dự đoán lối sống thực khuẩn thể. Như đã phân tích chi tiết trong Chương~\ref{chap:thuc_nghiem_danh_gia}, PhaBERT không chỉ cải thiện độ chính xác dự đoán trên contig ngắn và trung bình mà còn khắc phục được các hạn chế quan trọng của các phương pháp hiện có như sự phụ thuộc vào công cụ dự đoán gen và khả năng xử lý hạn chế trên các contig ngắn, phân mảnh. Sự kết hợp giữa biểu diễn ngữ cảnh toàn cục từ DNABERT-2 và trích xuất đặc trưng cục bộ chuyên biệt từ CNN đa cửa sổ tạo ra lợi thế rõ rệt so với từng thành phần riêng lẻ, mở ra hướng tiếp cận triển vọng cho việc tích hợp các mô hình tiền huấn luyện với các thành phần kiến trúc chuyên biệt cho các tác vụ đặc thù.

\section{Đóng góp của nghiên cứu}
\label{sec:contributions_c5}

Về mặt phương pháp, luận văn đề xuất PhaBERT, một kiến trúc học sâu lai kết hợp mô hình nền tảng hệ gen DNABERT-2 với mô-đun MKC chuyên biệt cho bài toán phân loại hệ gen thực khuẩn thể. Kiến trúc nhánh song song kép bao gồm nhánh CNN đa cửa sổ với ba kích thước (3, 5, 7) để trích xuất các đặc trưng cục bộ ở nhiều tỷ lệ không gian, và gộp trung bình có mặt nạ để tổng hợp biểu diễn toàn cục từ toàn bộ các biểu diễn token của DNABERT-2. Chiến lược huấn luyện hai giai đoạn với tốc độ học phân biệt, giai đoạn một đóng băng mạng nền để huấn luyện đầu phân loại, giai đoạn hai tinh chỉnh toàn bộ mô hình với tốc độ học thấp hơn cho mạng nền, cho phép tối ưu hóa hiệu quả cả biểu diễn tiền huấn luyện lẫn các thành phần chuyên biệt cho tác vụ. Phương pháp này khác biệt với các hướng tiếp cận trước đây bằng cách hoạt động trực tiếp trên chuỗi DNA thô thông qua mã hóa BPE, loại bỏ sự phụ thuộc vào công cụ phát hiện gen và cho phép xử lý hiệu quả các contig ngắn.

Về mặt kết quả thực nghiệm, luận văn cung cấp đánh giá trên bốn nhóm độ dài contig từ 100 đến 1.800 bp thông qua kiểm định chéo phân tầng 5 phần, cho thấy PhaBERT đạt độ chính xác cao nhất trên tất cả bốn nhóm với mức cải thiện 0,36--6,65\% so với các phương pháp tiên tiến. Nghiên cứu loại bỏ thành phần với bảy cấu hình gồm đường cơ sở DNABERT-2, PhaBERT, DNABERT-2-CNN (chỉ nhánh CNN đa cửa sổ), DNABERT-2-AP (gộp chú ý), PhaBERT-NO-K5, PhaBERT-NO-K3 và PhaBERT-NO-K7, chứng minh rằng mô-đun MKC đóng góp cải thiện độ chính xác từ 0,49\% đến 3,02\% tùy theo độ dài contig, với mức cải thiện độ đặc hiệu đặc biệt đáng kể (8,61\%) trên contig ngắn. So sánh các thành phần riêng lẻ cho thấy gộp chú ý mang lại lợi thế trên contig ngắn, trong khi CNN đa cửa sổ phát huy hiệu quả trên contig trung bình và dài; sự kết hợp cả hai nhánh trong PhaBERT mang lại hiệu suất cân bằng và ổn định nhất trên toàn phổ độ dài. Kết quả loại bỏ từng kích thước cửa sổ tích chập cho thấy các nhánh có mức đóng góp tương đương và thông tin chồng lấn đáng kể, xác nhận rằng giá trị của kiến trúc đa tỷ lệ nằm ở sự kết hợp tổng thể chứ không phải ở bất kỳ kích thước cửa sổ riêng lẻ nào. Kết quả này xác nhận giả thuyết rằng các mô hình nền tảng hệ gen cung cấp nền tảng biểu diễn mạnh mẽ cho các tác vụ hạ nguồn, nhưng vẫn cần các thành phần kiến trúc chuyên biệt để đạt hiệu suất tối ưu. Phân tích chi tiết về hiệu suất trên từng nhóm độ dài cũng cung cấp những nhận định quan trọng về điểm mạnh và hạn chế của từng phương pháp, giúp định hướng cho các nghiên cứu tiếp theo.

Về mặt ứng dụng thực tiễn, PhaBERT giải quyết một thách thức quan trọng trong phân tích dữ liệu hệ gen môi trường: phân loại hệ gen thực khuẩn thể từ các contig phân mảnh thu được qua quá trình lắp ráp hệ gen môi trường. Khả năng xử lý trực tiếp trên chuỗi DNA thô mà không cần công cụ dự đoán gen như Prodigal~\cite{hyatt2010prodigal} làm cho phương pháp đặc biệt phù hợp với các tình huống thực tế, nơi contig thường ngắn và thiếu các vùng mã hóa protein hoàn chỉnh. Hiệu suất cân bằng giữa độ nhạy và độ đặc hiệu trên tất cả các nhóm độ dài cho thấy PhaBERT có tiềm năng tích hợp vào các quy trình phân tích hệ gen môi trường tự động, hỗ trợ nghiên cứu về tương tác giữa thực khuẩn thể và vi khuẩn, động lực quần thể vi sinh vật, và ứng dụng liệu pháp thực khuẩn thể. Hơn nữa, khuôn khổ lai ghép giữa mô hình nền tảng tiền huấn luyện với các thành phần kiến trúc chuyên biệt cho tác vụ mở ra một hướng tiếp cận có thể mở rộng sang nhiều bài toán tin sinh học hạ nguồn khác như dự đoán vi khuẩn chủ, phát hiện prophage, phân loại chức năng gen và chú giải hệ gen vi sinh vật, qua đó góp phần thúc đẩy nghiên cứu trong lĩnh vực tin sinh học tính toán.

Tầm quan trọng của việc cải thiện độ chính xác phân loại được thể hiện rõ hơn khi đặt trong bối cảnh quy mô và hệ quả của dữ liệu hệ gen môi trường. Các cơ sở dữ liệu vi-rút từ hệ gen môi trường như IMG/VR phiên bản 4 hiện lưu trữ hơn 15 triệu hệ gen và mảnh hệ gen vi-rút~\cite{camargo2023imgvr}, cho thấy khối lượng trình tự cần phân loại trong các phân tích hạ nguồn là rất lớn. Ở quy mô này, ngay cả những sai lệch nhỏ về tỷ lệ phân loại cũng có thể tích lũy thành một số lượng đáng kể các trình tự bị gán nhãn sai. Hệ quả của việc phân loại sai lối sống không chỉ giới hạn ở khía cạnh thống kê: việc nhầm lẫn giữa thực khuẩn thể độc lực và ôn hòa có thể dẫn tới những suy luận sai lệch về động lực quần thể vi sinh vật, tiềm năng chuyển gen ngang, cũng như tính phù hợp của thực khuẩn thể cho ứng dụng liệu pháp~\cite{shang2023phatyp, wu2021deephage}. Do đó, mức cải thiện 0,36--6,65\% mà PhaBERT đạt được so với các phương pháp tiên tiến trên tập đánh giá của luận văn (Bảng~\ref{tab:main_results}) có thể mang lại lợi ích thực tiễn đáng kể khi được áp dụng trên các tập dữ liệu hệ gen môi trường quy mô lớn, mặc dù mức độ cụ thể còn phụ thuộc vào đặc điểm phân bố và chất lượng trình tự của từng tập dữ liệu.

\section{Hạn chế và hướng phát triển tương lai}
\label{sec:limitations_future}

Một hạn chế của PhaBERT là việc sử dụng lấy mẫu giảm ngẫu nhiên để xử lý mất cân bằng nhãn, vốn có thể làm mất một phần thông tin từ lớp đa số. Phân tích lỗi theo lớp nhãn (Phần~\ref{sec:loi_theo_lop}) cho thấy mô hình vẫn còn thiên lệch về phía lớp đa số. Các nghiên cứu tiếp theo có thể khảo sát những chiến lược giữ trọn dữ liệu huấn luyện mà vẫn hạn chế thiên lệch, chẳng hạn gán trọng số lớp hoặc sử dụng hàm Focal Loss~\cite{lin2017focal}.

Hạn chế thứ hai liên quan đến việc PhaBERT hoạt động thuần túy trên chuỗi nucleotide, chưa khai thác thông tin ở cấp độ protein. Mặc dù PhaBERT đạt độ chính xác cao nhất trên tất cả bốn nhóm contig, PhaTYP~\cite{shang2023phatyp} vẫn đạt độ đặc hiệu cao hơn trên contig dài (Nhóm D: 89,75\% so với 87,63\%) nhờ khai thác đặc trưng protein. Điều này cho thấy thông tin ở cấp độ protein cung cấp tín hiệu phân loại bổ sung quan trọng, đặc biệt cho việc nhận diện thực khuẩn thể ôn hòa trên các trình tự dài chứa đầy đủ các vùng mã hóa protein.

Cuối cùng, luận văn chưa đánh giá mô hình trên một tập dữ liệu độc lập hoàn toàn, cũng như chưa so sánh với các mô hình nền tảng hệ gen quy mô lớn hơn như Nucleotide Transformer~\cite{dalla2025nucleotide}; nghiên cứu mới chỉ sử dụng DNABERT-2 làm mạng nền, do chi phí tính toán của các mô hình lớn cao và nằm ngoài phạm vi tài nguyên hiện có. Hướng phát triển tiếp theo là kiểm chứng PhaBERT trên dữ liệu hệ gen môi trường độc lập và đánh giá mô-đun MKC trên nhiều mạng nền khác nhau, qua đó đánh giá chặt chẽ hơn khả năng tổng quát hóa của mô hình trong các điều kiện triển khai thực tế.